\subsection{Elektronenhülle und Periodensystem}
\label{subsec:elektronenhuelle-periodensystem}

Nach der klassischen Beschreibung des Elektrons im elektromagnetischen Feld
rückt nun eine neue Frage in den Mittelpunkt: Welche Rolle spielt das Elektron
in der Materie selbst? Die Antwort führt unmittelbar zur
Elektronenhülle\index{Elektronenhülle} der Atome und damit zum
Periodensystem\index{Periodensystem} der Elemente.

Ein Atom besteht aus einem positiv geladenen Kern und einer negativ geladenen
Elektronenhülle. Der Atomkern\index{Atomkern} enthält nahezu die gesamte Masse
des Atoms. Die räumliche Ausdehnung und die chemischen Eigenschaften eines
Atoms werden jedoch wesentlich durch die Elektronenhülle bestimmt.

Die Elektronen sind dabei nicht beliebig um den Kern angeordnet. Sie besetzen
bestimmte Zustände, die in einer ersten anschaulichen Näherung als
Schalen\index{Elektronenschale} beschrieben werden können. Jede Schale kann
nur eine begrenzte Anzahl von Elektronen aufnehmen. Dadurch entsteht eine
geordnete Struktur der Elektronenhülle.

Diese Ordnung ist der Schlüssel zum Periodensystem. Die Elemente sind dort
nach steigender Kernladungszahl\index{Kernladungszahl} angeordnet. Mit jedem
zusätzlichen Proton im Kern kommt bei einem neutralen Atom auch ein weiteres
Elektron in der Hülle hinzu. Die Art, wie diese Elektronen die verfügbaren
Zustände auffüllen, führt zur periodischen Wiederkehr ähnlicher chemischer
Eigenschaften.

Diese Ordnung wird im Periodensystem sichtbar
(Abbildung~\ref{fig:periodensystem-elektronenhuelle}).

\begin{figure}[htbp]
	\centering
	\includegraphics[width=\textwidth]{bilder/periodensystem.png}
	\caption{Das Periodensystem der Elemente als Ausdruck der Elektronenhülle.
		Die Elemente sind nach steigender Kernladungszahl angeordnet. Elemente in
		derselben Gruppe besitzen ähnliche äußere Elektronenstrukturen und zeigen
		deshalb oft ähnliche chemische Eigenschaften.}
	\label{fig:periodensystem-elektronenhuelle}
\end{figure}

\begin{PhysicsBox}[Elektronenhülle und Periodensystem]
	\label{box:elektronenhuelle-und-periodensystem}
	Das Periodensystem ist nicht nur eine tabellarische Ordnung chemischer
	Elemente. Es ist ein sichtbarer Ausdruck der Struktur der Elektronenhülle.
	
	Elemente mit ähnlicher äußerer Elektronenstruktur besitzen ähnliche chemische
	Eigenschaften. Deshalb stehen sie im Periodensystem in derselben Gruppe.
\end{PhysicsBox}

Besonders wichtig sind die Elektronen der äußersten besetzten Schale. Sie
bestimmen maßgeblich, wie ein Atom mit anderen Atomen wechselwirkt. Diese
äußeren Elektronen werden Valenzelektronen\index{Valenzelektronen} genannt.
Sie stehen im Mittelpunkt des nächsten Abschnitts.

Das Schalenmodell ist dabei nur eine vereinfachte Darstellung. In der modernen
Quantenmechanik beschreibt man Elektronen nicht als kleine Teilchen auf festen
Kreisbahnen, sondern durch Zustände und Wahrscheinlichkeitsverteilungen. Für
das Verständnis des Periodensystems bleibt die Schalenstruktur dennoch eine
sehr nützliche anschauliche Näherung.

\begin{DidacticBox}[Warum das Periodensystem periodisch ist]
	\label{box:warum-periodensystem-periodisch-ist}
	Die chemischen Eigenschaften der Elemente wiederholen sich nicht zufällig.
	Sie hängen eng mit dem Aufbau der Elektronenhülle zusammen.
	
	Wenn sich die äußere Elektronenstruktur in ähnlicher Weise wiederholt,
	treten auch ähnliche chemische Eigenschaften auf. Genau diese Wiederkehr
	macht das Periodensystem zu einer Ordnung der Materie.
\end{DidacticBox}

Damit zeigt sich ein zentraler Zusammenhang: Die Ordnung des Periodensystems
ist keine bloße chemische Sortierung. Sie beruht auf der inneren Struktur der
Elektronenhülle. Das Elektron ist daher nicht nur ein Träger elektrischer
Ladung, sondern der entscheidende Baustein für die chemische Vielfalt der
Materie.

Im nächsten Abschnitt wird diese Idee weitergeführt. Dort geht es um die
Valenzelektronen und ihre Rolle bei der chemischen Bindung.